
Reinforcement learning from execution feedback (RLEF) for code generation trains a language model to produce programs that pass test cases, using execution outcomes as reward signals. Within this framework, Group Relative Policy Optimization (GRPO) \citep{Shao2024} computes per-rollout advantages by normalizing rewards relative to other rollouts within the same prompt group: $A_i$ = (r\_i $-$ mean(r)) / std(r). This normalization eliminates the need for a learned value network, making GRPO computationally attractive for large language model fine-tuning.

A well-known failure mode of GRPO with binary execution rewards is variance collapse. When the reward function assigns 0 to any solution that does not pass all test cases, and when the model cannot produce a fully-passing solution, every rollout in a group receives reward 0. The group standard deviation is 0, and the advantages are undefined or zero. No gradient flows, and the model cannot learn from the problem. This failure mode is particularly acute on difficult problems, precisely where the model most needs to improve.

Ratio rewards --- assigning each rollout a score equal to the fraction of test cases passed --- address variance collapse by providing partial credit. Even when no rollout passes all tests, rollouts that pass a subset receive nonzero, differentiated rewards, preserving within-group variance. The theoretical question is how large this variance advantage is, and whether it is large enough to justify a modified training protocol.

This paper addresses three components of this question. First, it provides an analytical derivation of the expected within-group variance under both reward functions as a function of problem tractability q and test count T, and characterizes the variance ratio $\rho$(q, T) over the parameter space relevant to APPS introductory problems. Second, it formalizes a prescreening criterion --- the group tractability score S\_term --- and defines the partial-tractability window S\_term $\in$ [0.3, 0.55] where the variance advantage is predicted to be largest. Third, it reports the results of applying this prescreening protocol to 300 APPS introductory problems with Qwen2.5-Coder-7B-Instruct, including the finding that the base model achieves 0\% pass rate across 2,400 solution attempts, blocking all subsequent mechanism experiments.

The paper makes four contributions:

\begin{enumerate}
\item A formal derivation of the Binomial variance advantage of ratio rewards over binary rewards under GRPO's group-relative normalization, with analytical characterization of the variance ratio as a function of q and T.
\item The formalization of S\_term as a computable group-level tractability score and the definition of the partial-tractability window S\_term $\in$ [0.3, 0.55] through this analytical framework.
\item A production-validated prescreening pipeline (15/15 tasks, 67/67 integration tests) implementing the prescreening protocol for APPS introductory problems.
\item An empirical finding that Qwen2.5-Coder-7B-Instruct achieves 0\% pass rate on APPS introductory problems under a strict execution harness without SFT initialization, with an analysis of the probable contributing factors.
\end{enumerate}

The mechanism hypotheses --- whether ratio rewards produce more distinct advantage levels, higher gradient covariance, higher gradient SNR, and earlier zero-reward-fraction escape than binary rewards --- could not be evaluated because no problems qualified for prescreening. These hypotheses remain in a blocked state pending an SFT-initialized model checkpoint.

